% !TeX root = ../tesis.tex
%%%%%%%%%%%%%%%%%%%%%%%%%%%%%%%%%%%%%%%%%%%%%%%%%%%%%%%%%%%%%%%%%%%
% SECCIÓN 5.2-2 — RESULTADOS VFT: FASE 2
% Ponderación dinámica de la red
%%%%%%%%%%%%%%%%%%%%%%%%%%%%%%%%%%%%%%%%%%%%%%%%%%%%%%%%%%%%%%%%%%%
\section{\bajoRevision{Fase 2: Ponderación Dinámica de la Red}}
\label{sec:resultados-fase2}

\notaRevision{Los indicadores de esta fase se encuentran en revisión con el
comité tutor. Los coeficientes de fricción vial y la penalización por
transferencia han sido calculados e integrados al grafo; el \termino{Node Score}
se encuentra en etapa de planificación, pendiente de la disponibilidad de las
vistas materializadas MV1 y MV2 de \textbf{Apimetro}.}

La Fase~2 del Modelo VFT enriquece el grafo construido en la Fase~1 con dos
capas de pesos dinámicos que aproximan las condiciones reales de operación de la
red: el \textbf{Coeficiente de Fricción Vial} ($CF$), que penaliza el tiempo de
traslado en función del nivel de interferencia del tráfico urbano sobre cada
tipo de vía, y la \textbf{Penalización por Transferencia} ($T_b$), que incorpora
el costo de caminata y espera en cada transbordo intermodal. Estos dos pesos
conforman el componente dinámico de la \textbf{Ecuación VFT}:

\begin{equation}
    w(u,v) = \frac{D_{uv}}{V_{\text{modo}}} \cdot CF_{\text{vía}} + T_b
    \label{eq:vft-arista}
\end{equation}

\begin{itemize}
    \item $D_{uv}$: distancia geodésica del segmento calculada mediante la
    fórmula de \termino{Haversine} (metros).
    \item $V_{\text{modo}}$: velocidad comercial ajustada del sistema de
    transporte (km/h), extraída del \geojson{} de \textbf{Apimetro} o
    sustituida por el valor de \termino{fallback} calibrado por sistema.
    \item $CF_{\text{vía}}$: coeficiente de fricción vial según el tipo de
    derecho de vía del segmento (adimensional, $\geq 1$).
    \item $T_b$: penalización por transferencia peatonal en minutos; se aplica
    exclusivamente en las aristas de transbordo generadas por el
    \termino{Snapping Lógico Geográfico}.
\end{itemize}

Para las aristas de servicio normal, el peso de enrutamiento es
$\mathit{weight} = (D_{uv}/V_{\text{modo}}) \cdot CF_{\text{vía}}$; el
$T_b$ queda registrado como atributo informativo pero no penaliza el
segmento. La penalización por transferencia actúa únicamente sobre aristas de
tipo \texttt{transfer}, como se describe en la Sección~\ref{subsec:res-penalizacion-transferencia}.

% ----------------------------------------------------------------
\subsection{Coeficiente de Fricción Vial ($CF$)}
\label{subsec:coeficiente-friccion}
% ----------------------------------------------------------------

El \textbf{Coeficiente de Fricción Vial} ($CF$) traduce el tipo de
\termino{derecho de vía} de cada segmento en un penalizador multiplicativo
sobre el tiempo de traslado. La premisa es que circular por un túnel del
Metro —aislado del tráfico en superficie— implica un costo temporal muy
distinto al de un camión de la \textsc{rtp} inmerso en el flujo vehicular de
Insurgentes o Periférico.

\subsubsection{Constante de Saturación Vial}
\label{subsubsec:beta-saturacion}

El parámetro base del coeficiente es la constante de saturación $\beta$,
derivada del \termino{TomTom Traffic Index} para la \zmvm:

\begin{equation}
    \beta_{\text{ZMVM}} = 0.759
    \label{eq:beta-zmvm}
\end{equation}

Este valor indica que la ciudad opera, en condiciones promedio, al 76~\% de
su capacidad vial. Es decir, un vehículo en tráfico mixto experimenta, en
promedio, un sobretiempo del 76~\% respecto al recorrido en flujo libre.

\subsubsection{Fórmula y Valores por Derecho de Vía}
\label{subsubsec:formula-cf}

El coeficiente $CF$ se obtiene a partir de un parámetro de escala $\alpha$
que modula la exposición de cada tipo de vía a la saturación urbana:

\begin{equation}
    CF = 1 + \alpha \cdot \beta_{\text{ZMVM}}
    \label{eq:cf-formula}
\end{equation}

\begin{itemize}
    \item $\alpha \in \{0.0,\ 0.2,\ 0.5,\ 1.0\}$: fracción de la saturación
    urbana que afecta el segmento, determinada por su tipo de derecho de vía.
    \item $\beta_{\text{ZMVM}} = 0.759$: constante de saturación de la \zmvm
    (\termino{TomTom Traffic Index}).
    \item $CF \geq 1$: un valor de 1.0 indica flujo ideal sin interferencia;
    valores superiores representan el penalizador proporcional al congestionamiento.
\end{itemize}

El Cuadro~\ref{tab:cf-por-via} sintetiza los cuatro tipos de derecho de vía
modelados, los sistemas de transporte correspondientes y el valor numérico del
coeficiente resultante.

\begin{table}[H]
    \centering
    \caption[Coeficiente de Fricción Vial por tipo de derecho de vía]{Coeficiente
    de Fricción Vial ($CF$) por tipo de derecho de vía en el Modelo VFT.
    $\alpha$ escala el impacto de la saturación urbana
    $\beta_{\text{ZMVM}} = 0.759$.}
    \label{tab:cf-por-via}
    \begin{tabularx}{\textwidth}{|l|c|c|X|}
        \hline
        \textbf{Derecho de vía} & \textbf{$\alpha$} & \textbf{$CF$} &
        \textbf{Sistemas representativos} \\
        \hline
        Exclusivo   & 0.0 & 1.000 & Metro, Tren Suburbano \\
        Confinado   & 0.2 & 1.152 & Metrobús, Cablebús, Mexicable \\
        Compartido  & 0.5 & 1.380 & Trolebús, Tren Ligero \\
        Mixto       & 1.0 & 1.759 & \textsc{rtp}, Corredores Concesionados,
                                    Mexibús, Pumabús \\
        \hline
    \end{tabularx}
    \fuentefigura{Elaboración propia con base en \texttt{impedance.py}
    (VFTModel, corrida 2026-04-18). Los valores $\alpha$ tienen como referencia
    normativa el \textit{Manual de Calles: Diseño Vial para Ciudades Mexicanas}
    \parencite{sedatu2019manual}.}
\end{table}

El rango del coeficiente $CF \in [1.000,\ 1.759]$ implica que el tiempo de
traslado en un segmento de vía mixta (RTP, Corredores Concesionados) puede ser
hasta un 75.9~\% mayor que el mismo recorrido bajo condiciones de flujo ideal.
Esta diferencia no refleja únicamente la velocidad operacional del vehículo —que
ya está calibrada en $V_{\text{modo}}$— sino el efecto sistemático de las
detenciones por semáforos, cruces peatonales y congestión que son inherentes a
los segmentos de superficie compartida.

En el caso del sistema de derecho de vía desconocido, el Modelo VFT asigna
por defecto el valor de vía \texttt{mixto} ($CF = 1.759$), aplicando el
principio de escenario más conservador para no subestimar los tiempos de traslado.

% ----------------------------------------------------------------
\subsection{Penalización por Transferencia ($T_b$)}
\label{subsec:res-penalizacion-transferencia}
% ----------------------------------------------------------------

La \textbf{Penalización por Transferencia} ($T_b$) modela el costo temporal que
un pasajero incurre al cambiar de sistema o línea dentro de la red. A diferencia
del tiempo de traslado —que depende de la distancia y la velocidad del vehículo—,
el costo de transbordo consta de dos componentes independientes: el tiempo de
caminata entre plataformas y el tiempo de espera en el andén de destino.

\subsubsection{Fórmula de la Penalización}
\label{subsubsec:formula-tb}

Para cada arista de transbordo generada por el \termino{Snapping Lógico Geográfico}
entre los nodos $u$ (origen del transbordo) y $v$ (destino del transbordo), el
peso de la arista es:

\begin{equation}
    T_b(u, v) = t_{\text{caminata}}(u, v) + \frac{F_v}{2}
    \label{eq:tb-formula}
\end{equation}

\begin{itemize}
    \item $t_{\text{caminata}}(u, v)$: tiempo de caminata entre los dos nodos
    a una velocidad peatonal normativa de 5~km/h (1.38~m/s), equivalente al
    límite superior establecido en el \textit{Manual de Calles SEDATU}
    \parencite{sedatu2019manual} para infraestructura de transbordo sin detenciones
    en semáforo.
    \item $F_v$: frecuencia de servicio del sistema destino $v$ (minutos entre
    unidades); si no está reportada en el \geojson{}, se utiliza el valor
    de \termino{fallback} calibrado por sistema.
    \item $F_v / 2$: tiempo de espera esperado, calculado como la mitad de la
    frecuencia bajo la hipótesis de llegada aleatoria del pasajero al andén
    (\termino{random arrival assumption}).
\end{itemize}

El tiempo de caminata se calcula como:

\begin{equation}
    t_{\text{caminata}}(u, v) = \frac{D_{uv}}{V_{\text{peatonal}}}
    = \frac{D_{uv}}{83.33\ \text{m/min}}
    \label{eq:t-caminata}
\end{equation}

donde $D_{uv}$ es la distancia geodésica entre los dos nodos (Haversine) en
metros, y $V_{\text{peatonal}} = 5000\ \text{m} / 60\ \text{min} = 83.33\ \text{m/min}$.

\subsubsection{Valores de Frecuencia de Fallback}
\label{subsubsec:fallback-frecuencia}

El Cuadro~\ref{tab:fallback-frecuencia} presenta los valores de frecuencia de
\termino{fallback} calibrados para cada sistema de transporte, junto con el
costo de abordaje resultante ($F_v / 2$). Estos valores representan intervalos
promedio en condiciones de operación normal en hora no pico.

\begin{table}[H]
    \centering
    \caption[Frecuencias de \textit{fallback} y costo de abordaje por sistema]{Frecuencias
    de servicio de \termino{fallback} y penalización de abordaje resultante
    ($T_b^{\text{espera}} = F_v / 2$) por sistema de transporte en el Modelo VFT.}
    \label{tab:fallback-frecuencia}
    \begin{tabular}{|l|c|c|}
        \hline
        \textbf{Sistema} & \textbf{$F_v$ (min)} & \textbf{$F_v / 2$ (min)} \\
        \hline
        \textsc{Metro} (STC)          & 3.0  & 1.50 \\
        Metrobús (MB)                 & 5.0  & 2.50 \\
        Tren Ligero (TL)              & 10.0 & 5.00 \\
        Trolebús (TROLE)              & 8.0  & 4.00 \\
        \textsc{rtp}                  & 15.0 & 7.50 \\
        Tren Suburbano (SUB)          & 12.0 & 6.00 \\
        Cablebús (CBB)                & 10.0 & 5.00 \\
        Corredores Concesionados (CC) & 15.0 & 7.50 \\
        Mexibús                       & 6.0  & 3.00 \\
        \hline
    \end{tabular}
    \fuentefigura{Elaboración propia con base en \texttt{graph\_builder.py}
    (VFTModel, corrida 2026-04-18). Los valores corresponden a frecuencias
    promedio en horario valle; las frecuencias en hora pico son reportadas por
    los operadores como atributo del \geojson{} y tienen precedencia sobre los
    \termino{fallbacks}.}
\end{table}

\subsubsection{Decisión de Diseño: Boarding Cost Informativo}
\label{subsubsec:boarding-cost-diseno}

Una precisión metodológica relevante: el costo de abordaje $F_v / 2$ \textbf{no
se suma al peso de las aristas de servicio normal}. Solo actúa como componente
del peso de las aristas de transbordo peatonal. Esta decisión responde a que el
pasajero ya está dentro del vehículo en los segmentos normales; el tiempo de
espera por frecuencia solo es pertinente en el momento de abordar un nuevo
sistema. Sumarlo en cada arista normal equivaldría a penalizar al pasajero en
cada estación intermedia, sobreestimando sistemáticamente los tiempos de ruta.

En consecuencia, el Modelo VFT registra $F_v / 2$ como atributo informativo
en las aristas de servicio (\texttt{boarding\_cost\_min}) para análisis
descriptivos, pero el atributo \texttt{weight} de estas aristas solo contiene
el tiempo de traslado con fricción.

% ----------------------------------------------------------------
\subsection{\bajoRevision{Node Score — Estado de Implementación}}
\label{subsec:node-score}
% ----------------------------------------------------------------

\notaRevision{El \termino{Node Score} ($\text{VFT\_score}$) se encuentra en etapa
de planificación. Su cálculo requiere los datos de afluencia por línea (MV1) y
los polígonos AGEB con atributos censales (MV2) de \textbf{Apimetro}, aún no
disponibles en la fecha de cierre de este documento. La arquitectura del indicador,
sus fórmulas y las dependencias de datos se presentan a continuación como
fundamento para la etapa de implementación.}

El \termino{Node Score} ($\text{VFT\_score}_i$) constituye la segunda capa de
pesos del Modelo VFT: mientras que los pesos de arista modelan el costo temporal
de recorrer la red, el \termino{Node Score} asigna a cada nodo una medida de
\textbf{importancia compuesta} que combina propiedades topológicas de la red con
información de demanda observada y potencial. Este indicador es el que alimenta
el análisis de vulnerabilidad de la Fase~3.

\subsubsection{Fórmula del Node Score}
\label{subsubsec:formula-vft-score}

\begin{equation}
    \text{VFT\_score}_i = \alpha_1 \hat{k}_i + \alpha_2 \hat{a}_i + \alpha_3 \hat{p}_i
    \label{eq:vft-score}
\end{equation}

donde los tres componentes normalizados al rango $[0, 1]$ son:

\begin{itemize}
    \item $\hat{k}_i$: \textbf{Fuerza Capilar normalizada} (min-max sobre todos
    los nodos del grafo). Calculada en la Fase~1 para los 11{,}115 nodos.
    \item $\hat{a}_i$: \textbf{Afluencia estimada por estación normalizada}.
    Derivada de la afluencia diaria por línea ($A_\ell$, fuente MV1) desagregada
    internamente por VFT mediante ponderación por grado topológico:
    \begin{equation}
        a_i = \sum_{\ell \ni i} A_\ell \cdot
        \frac{g_i}{\displaystyle\sum_{j \in \ell} g_j}
        \label{eq:desagregacion-afluencia}
    \end{equation}
    donde $g_i = \deg(i)$ es el grado del nodo $i$ en el grafo NetworkX.
    \item $\hat{p}_i$: \textbf{Población captada en buffer de 800~m normalizada}.
    Calculada intersectando un radio de 800~m alrededor de cada estación con los
    polígonos AGEB (fuente MV2), ponderando cada AGEB por la fracción de su área
    dentro del buffer:
    \begin{equation}
        p_i = \sum_{j \in \mathcal{A}_i} \text{pob}_j \cdot
        \frac{\text{área}(B_i \cap \text{AGEB}_j)}{\text{área}(\text{AGEB}_j)}
        \label{eq:poblacion-captada}
    \end{equation}
    donde $B_i$ es el buffer circular de 800~m en proyección métrica EPSG:6372
    (México Centro) y $\mathcal{A}_i$ el conjunto de AGEBs que lo intersectan.
\end{itemize}

La normalización de cada componente usa el método \termino{min-max}:
\begin{equation}
    \hat{x}_i = \frac{x_i - \min(\mathbf{x})}{\max(\mathbf{x}) - \min(\mathbf{x})}
    \label{eq:normalizacion-minmax}
\end{equation}

\subsubsection{Determinación de los Pesos $\alpha$}
\label{subsubsec:pesos-alpha}

Los pesos $\alpha_1, \alpha_2, \alpha_3$ se determinarán en dos etapas
progresivas una vez que los datos de MV1 y MV2 estén disponibles:

\begin{enumerate}
    \item \textbf{Validación inicial (pesos iguales):} $\alpha_1 = \alpha_2 =
    \alpha_3 = 1/3$. Objetivo: verificar que nodos de alta importancia reconocida
    (Pantitlán, Insurgentes, Pino Suárez, Balderas) aparecen en el top-10 del
    \termino{ranking}. Si no lo hacen, el problema se ubica en los datos de entrada.
    \item \textbf{Pesos por entropía (CRITIC):} los pesos finales se derivan de
    la variabilidad observada de cada componente sobre todos los nodos:
    \begin{equation}
        \alpha_n = \frac{\sigma_n}{\displaystyle\sum_{m=1}^{3} \sigma_m}
        \label{eq:alpha-entropia}
    \end{equation}
    donde $\sigma_n = \text{std}(\hat{x}_n)$. Un componente más discriminante
    (mayor dispersión entre nodos) recibe mayor peso, penalizando automáticamente
    los componentes redundantes.
\end{enumerate}

\subsubsection{Dependencias de Datos y Estado}
\label{subsubsec:node-score-estado}

El Cuadro~\ref{tab:node-score-dependencias} resume los componentes del
\termino{Node Score}, su fuente de datos y su estado de disponibilidad a la
fecha de cierre del documento.

\begin{table}[H]
    \centering
    \caption[Dependencias del Node Score por componente]{Componentes del
    $\text{VFT\_score}$, fuentes de datos requeridas y estado de implementación.}
    \label{tab:node-score-dependencias}
    \begin{tabularx}{\textwidth}{|l|X|l|l|}
        \hline
        \textbf{Componente} & \textbf{Fuente} & \textbf{Responsable} & \textbf{Estado} \\
        \hline
        $\hat{k}_i$ Fuerza Capilar    & VFT interno (\texttt{capillary\_strength.py})
                                        & VFTModel & Calculado \\
        $\hat{a}_i$ Afluencia         & MV1 Apimetro — afluencia por línea
                                        & Apimetro & Pendiente \\
        $\hat{p}_i$ Población captada & MV2 Apimetro — AGEBs Censo INEGI 2020
                                        & Apimetro & Pendiente \\
        Pesos $\alpha$ (entropía)     & Derivados de $\sigma(\hat{k}, \hat{a}, \hat{p})$
                                        & VFTModel & Pendiente MV1+MV2 \\
        \hline
    \end{tabularx}
    \fuentefigura{Elaboración propia con base en
    \texttt{weighted\_graph\_enrichment.md} y
    \texttt{apimetro\_requerimientos\_mvs.md} (VFTModel, 2026-04-28).}
\end{table}

El diseño del \termino{Node Score} garantiza la \textbf{independencia entre
componentes}: la desagregación de afluencia usa el grado topológico crudo
$g_i = \deg(i)$ —no la Fuerza Capilar $k_i$— para evitar que $\hat{k}_i$
domine doblemente el cómputo. Antes de asignar los pesos finales por entropía,
se verificará la correlación $\rho(\hat{k}, \hat{a})$: si $\rho \geq 0.75$,
los dos componentes se fusionarán en un índice combinado $\delta_i = \beta
\hat{k}_i + (1-\beta)\hat{a}_i$ antes de calcular el \termino{score} final.
